\begin{theorem}[Stokes 转迁的显式相对上同调模型与可计算基]\label{thm:cdim-stokes-explicit-relative-cohomology-basis}
\leanverified{paper\_cdim\_stokes\_explicit\_relative\_cohomology\_basis}
设 \(v\ge 1\)，并将闭单位圆盘 \(\overline{\mathbb D}\) 视为带角流形。
在 \(\overline{\mathbb D}^{\,v}\) 上取极坐标 \(z_j=r_j e^{\mathrm{i}\phi_j}\)（\(1\le j\le v\)）。
对任意非空指标集 \(I\subseteq[v]\) 定义微分形式
\[
\alpha_{I}:=\bigwedge_{j\in I}\left(\frac{r_j^2}{2\pi}\,d\phi_j\right)\in\Omega^{|I|}(\overline{\mathbb D}^{\,v}),
\qquad
\Omega_{I}:=d\alpha_I\in\Omega^{|I|+1}(\overline{\mathbb D}^{\,v}).
\]
则：
\begin{enumerate}
  \item \(\Omega_I|_{\TT^{v}}=0\)，从而 \(\Omega_I\) 定义一个相对闭类
  \(
  [\Omega_I]\in H^{|I|+1}(\overline{\mathbb D}^{\,v},\TT^{v};\ZZ)
  \)；
  并且当 \(|I|=k\) 时，\(\{[\Omega_I]:|I|=k\}\) 构成 \(H^{k+1}(\overline{\mathbb D}^{\,v},\TT^{v};\ZZ)\) 的一组 \(\ZZ\)-基。
  \item 设 \(\delta:H^{k}(\TT^{v};\ZZ)\to H^{k+1}(\overline{\mathbb D}^{\,v},\TT^{v};\ZZ)\) 为配对长正合列的连接同态。
  则在 \(\TT^{v}\) 的标准外代数基
  \(
  \bigwedge_{j\in I}\frac{d\phi_j}{2\pi}\in H^{k}(\TT^{v};\ZZ)
  \)
  上有显式公式
  \[
  \boxed{\;
  \delta\!\left(\bigwedge_{j\in I}\frac{d\phi_j}{2\pi}\right)=[\Omega_I]\; } .
  \]
  特别地，\(\delta:H^{1}(\TT^{v};\ZZ)\xrightarrow{\sim}H^{2}(\overline{\mathbb D}^{\,v},\TT^{v};\ZZ)\) 为同构，
  且与命题 \ref{prop:cdim-stokes-exact-sequence-dictionary} 的坐标化描述一致。
  \item 对任意 \(u\ge 0\)，在识别
  \(
  (\mathfrak X_0(\ZZ^{u}\oplus\NN^{v}),\Sigma_0(\ZZ^{u}\oplus\NN^{v}))
  \cong
  (\TT^{u}\times \overline{\mathbb D}^{\,v},\TT^{u}\times\TT^{v})
  \)
  下，右乘 cup-product 与 Künneth 分解给出
  \[
  H^{m}\bigl(\mathfrak X_0,\Sigma_0;\ZZ\bigr)
  \cong
  \bigoplus_{p+q=m}H^{p}(\TT^{u};\ZZ)\otimes H^{q}(\overline{\mathbb D}^{\,v},\TT^{v};\ZZ),
  \]
  从而由 \(\TT^u\) 的标准外代数基 \([\eta]\) 与 \([\Omega_I]\) 的 cup-product
  \(
  [\eta]\smile[\Omega_I]
  \)
  得到一组显式可计算基；该基与命题 \ref{prop:cdim-stokes-relative-cohomology-ring} 的分次环分解相容。
\end{enumerate}
\end{theorem}
\begin{proof}
对任意 \(I\)，由于在边界 \(\TT^{v}\subset \overline{\mathbb D}^{\,v}\) 上各 \(r_j\equiv 1\) 为常数，故 \(dr_j|_{\TT^{v}}=0\)，从而 \(\Omega_I=d\alpha_I\) 限制到边界为零。
因此 \(\Omega_I\) 代表相对闭类。
\par
另一方面，\(\alpha_I|_{\TT^{v}}=\bigwedge_{j\in I}\frac{d\phi_j}{2\pi}\)。
由相对 de Rham 模型中连接同态的链级刻画，若 \(\beta\in\Omega^{k}(\overline{\mathbb D}^{\,v})\) 满足 \(\beta|_{\TT^v}\) 代表某个边界上同调类，则 \(\delta([\beta|_{\TT^v}])=[d\beta]\)。
取 \(\beta=\alpha_I\) 即得 \(\delta(\bigwedge_{j\in I}\frac{d\phi_j}{2\pi})=[\Omega_I]\)。
\par
由于 \(\overline{\mathbb D}^{\,v}\) 可缩，长正合列给出对 \(k\ge 1\) 的同构
\(
H^{k+1}(\overline{\mathbb D}^{\,v},\TT^{v};\ZZ)\cong H^{k}(\TT^{v};\ZZ)
\)，且由 \(\delta\) 实现；结合上式可知 \(\{[\Omega_I]:|I|=k\}\) 与 \(\TT^{v}\) 的标准外代数基一一对应，故构成 \(\ZZ\)-基。
\par
最后一条由配对 Künneth 公式与 \([\Omega_I]\) 的基性质直接推出，并与命题 \ref{prop:cdim-stokes-relative-cohomology-ring} 的环同构相容。证毕。
\end{proof}

\begin{theorem}[相对第一陈类与边界框架分类的 Stokes 同构]\label{thm:cdim-stokes-relative-c1-framed-line-bundles}
令 \(n\ge 1\)。设 \(\mathcal L^{\mathrm{fr}}(\overline{\mathbb D}^{\,n})\) 为带边界框架的复线丛同构类集合：
其元素为对 \((L,\tau)\)，其中 \(L\to \overline{\mathbb D}^{\,n}\) 为复线丛，\(\tau\) 为其在边界 \(\partial\overline{\mathbb D}^{\,n}=\TT^{n}\) 上的一个平凡化。
则存在自然同构
\[
\boxed{\;
c_1^{\mathrm{rel}}:\ \mathcal L^{\mathrm{fr}}(\overline{\mathbb D}^{\,n})\xrightarrow{\ \sim\ }H^{2}(\overline{\mathbb D}^{\,n},\TT^{n};\ZZ)\; } .
\]
并且当 \(L\) 为平凡丛 \(L=\overline{\mathbb D}^{\,n}\times\CC\) 且框架 \(\tau\) 由夹持函数 \(g_\tau:\TT^{n}\to U(1)\) 给出时，有
\[
\boxed{\;
c_1^{\mathrm{rel}}(L,\tau)=\delta([g_\tau])\in H^{2}(\overline{\mathbb D}^{\,n},\TT^{n};\ZZ)\; } ,
\]
其中 \(\delta:H^{1}(\TT^{n};\ZZ)\to H^{2}(\overline{\mathbb D}^{\,n},\TT^{n};\ZZ)\) 为连接同态（见定理 \ref{thm:cdim-stokes-explicit-relative-cohomology-basis} 与命题 \ref{prop:cdim-stokes-exact-sequence-dictionary}）。
此外，\(c_1^{\mathrm{rel}}\) 可由 Stokes 配对刻画为“边界绕数 \(=\) 体内曲率积分”：若取任意 unitary 连接 \(A\) 使其边界规范势与 \(\tau\) 相容，令 \(F=dA\)，则对任意相对 \(2\)-链 \(C\) 有
\[
\int_{C}F=\int_{\partial C}A.
\]
\leanverified{paper\_cdim\_stokes\_relative\_c1\_framed\_line\_bundles}
\end{theorem}
\begin{proof}
\(\overline{\mathbb D}^{\,n}\) 可缩，故任意复线丛在体内平凡；带边界框架的同构类因此等价于边界过渡函数的同伦类
\(
[\TT^{n},U(1)]\cong H^{1}(\TT^{n};\ZZ)
\)。
另一方面，由连接同态 \(\delta:H^{1}(\TT^{n};\ZZ)\xrightarrow{\sim}H^{2}(\overline{\mathbb D}^{\,n},\TT^{n};\ZZ)\)（定理 \ref{thm:cdim-stokes-explicit-relative-cohomology-basis}）得到自然同构，定义为 \(c_1^{\mathrm{rel}}\)。
\par
当 \(L\) 为平凡丛时，\(\tau\) 的同伦类由 \(g_\tau\) 给出，且上同构在该表示下即为 \(\delta([g_\tau])\)。
最后的 Stokes 配对刻画由 \(F=dA\) 与带边界的 Stokes 公式（命题 \ref{prop:cdim-stokes-half-face-count} 的局部模型）直接推出。证毕。
\end{proof}

\begin{proposition}[\texorpdfstring{$q$}{q}-幂映射的 Stokes 缩放律与异常放大的上同调实现]\label{prop:cdim-stokes-qpower-scaling-anom-amplification}
令 \(n\ge 1\)，并定义坐标 \texorpdfstring{$q$}{q}-幂映射
\[
P_q:\overline{\mathbb D}^{\,n}\to \overline{\mathbb D}^{\,n},
\qquad
(z_1,\dots,z_n)\mapsto (z_1^{q},\dots,z_n^{q})
\qquad(q\in\NN_{\ge 1}).
\]
则：
\begin{enumerate}
  \item 在相对上同调上，拉回
  \(
  P_q^\ast:H^{2}(\overline{\mathbb D}^{\,n},\TT^{n};\ZZ)\to H^{2}(\overline{\mathbb D}^{\,n},\TT^{n};\ZZ)
  \)
  等于标量乘法 \(q\)；
  更一般地，对任意非空 \(I\subseteq[n]\) 有
  \[
  \boxed{\;P_q^\ast[\Omega_I]=q^{|I|}\,[\Omega_I]\;}
  \quad\text{于}\quad
  H^{|I|+1}(\overline{\mathbb D}^{\,n},\TT^{n};\ZZ),
  \]
  其中 \([\Omega_I]\) 为定理 \ref{thm:cdim-stokes-explicit-relative-cohomology-basis} 的基类。
  \item 设 \(G\) 为有限阿贝尔群，并记异常签名向量空间
  \(
  \mathcal{A}_G:=\RR^{\widehat G\setminus\{\chi_0\}}
  \)
  （定义 \ref{def:pom-anom-accumulation}）。
  令 \(M_G:=\NN^{\widehat G\setminus\{\chi_0\}}\) 为自由正交交换幺半群，则
  \[
  (\mathfrak X_0(M_G),\Sigma_0(M_G))\cong (\overline{\mathbb D}^{\,n},\TT^{n}),\qquad n:=|\widehat G|-1,
  \]
  并存在一个按坐标基识别的自然线性同构
  \[
  \boxed{\;
  \mathcal{A}_G\ \cong\ H^{2}\bigl(\mathfrak X_0(M_G),\Sigma_0(M_G);\RR\bigr)\;}
  \]
  使得命题 \ref{prop:pom-anom-gap-amplification} 的放大律
  \(
  [\mathrm{Anom}(w^{\boxtimes q})]=q[\mathrm{Anom}(w)]
  \)
  与 \(P_q^\ast\) 的上同调缩放在该模型下对齐。
\end{enumerate}
\end{proposition}
\begin{proof}
(1) 在边界 \(\TT^{n}\) 上，\(P_q|_{\TT^{n}}\) 在每个圆周因子上为度数 \(q\) 的覆盖，故在 \(H^{1}(\TT^{n};\ZZ)\) 上为乘以 \(q\)，并在 \(\wedge^{|I|}\) 上为乘以 \(q^{|I|}\)。
由连接同态 \(\delta\) 的自然性与定理 \ref{thm:cdim-stokes-explicit-relative-cohomology-basis}(2) 的识别
\(
[\Omega_I]=\delta(\bigwedge_{j\in I}\frac{d\phi_j}{2\pi})
\)
立得 \(P_q^\ast[\Omega_I]=q^{|I|}[\Omega_I]\)；
取 \(|I|=1\) 即得相对 \(H^2\) 上的标量乘法 \(q\)。
\par
(2) 对 \(M_G=\NN^{n}\)，由定义 \ref{def:cdim-stokes-character-pair} 有
\(
\mathfrak X_0(M_G)\cong \overline{\mathbb D}^{\,n}
\)、
\(
\Sigma_0(M_G)\cong \TT^{n}
\)。
而 \(H^{2}(\overline{\mathbb D}^{\,n},\TT^{n};\RR)\cong \RR^{n}\) 的坐标基可取为 \(\{[\Omega_{\{j\}}]\}\)；
与 \(\mathcal{A}_G=\RR^{n}\) 的标准基对齐即可得到自然同构。
在该识别下，“放大律为乘以 \(q\)”与第 (1) 条在相对 \(H^2\) 上的缩放一致。证毕。
\end{proof}
\leanverified{paper\_cdim\_stokes\_qpower\_scaling\_anom\_amplification}

\endinput
